\documentclass[11pt,a4paper]{article}

\usepackage[margin=2.2cm]{geometry}
\usepackage{booktabs}
\usepackage{tabularx}
\usepackage{array}
\usepackage{enumitem}
\usepackage[hidelinks]{hyperref}
\usepackage{parskip}
\usepackage{url}
\urlstyle{tt}
\emergencystretch=1.5em

\setlist{nosep}
\newcommand{\partno}[1]{\texttt{#1}}
\newcolumntype{L}{>{\raggedright\arraybackslash}X}

\begin{document}

\begin{center}
{\LARGE\bfseries SecurePouch PCB Electrical Review Specification}\\[0.4em]
{\large nRF52840 + nRF9151 + nPM1300 smart pouch PCB}\\[0.8em]
\begin{tabularx}{0.92\textwidth}{lL}
\textbf{Status} & Electrical review draft\\
\textbf{Last updated} & June 2026\\
\textbf{Primary schematic} & \texttt{smart\_pouch.kicad\_sch}\\
\textbf{Purpose} & Reviewer checklist for power, USB-C, battery, PMIC, and high-current loads\\
\end{tabularx}
\end{center}

\section{Review Scope}

This document is intentionally limited to electrical design information needed for schematic review.
Mechanical, backend, app, and product-roadmap details have been removed from this review copy.

\section{System Summary}

\begin{table}[h]
\small
\begin{tabularx}{\textwidth}{llL}
\toprule
\textbf{Block} & \textbf{Schematic ref.} & \textbf{Notes}\\
\midrule
BLE/application MCU & U1, nRF52840 & Local peripherals, BLE, USB data, PMIC control over I2C/TWI\\
Cellular/GNSS SiP & U2, nRF9151-LACA-R7 & LTE-M/NB-IoT/GNSS; UART link to nRF52840\\
PMIC/charger & U?, nPM1300-CAAA-R7 & USB-C input, battery charging, BUCK1/BUCK2, fuel/ADC monitoring\\
Battery connector & CN1 & Intended 3-pin, 2.54 mm battery connector; pinout: VBAT+, NTC, GND\\
USB-C receptacle & USB1, C2765186 & 16-pin USB2.0 Type-C receptacle\\
Accelerometer & BMA400 & I2C, interrupt lines to nRF52840\\
Siren boost & MT3608 & 15 V rail for piezo H-bridge\\
Solenoid boost & MT3608 & 12 V rail for lock solenoid pulse\\
\bottomrule
\end{tabularx}
\end{table}

\section{Power Architecture}

\begin{table}[h]
\small
\begin{tabularx}{\textwidth}{lllL}
\toprule
\textbf{Rail/net} & \textbf{Source} & \textbf{Nominal} & \textbf{Consumers / review notes}\\
\midrule
VBAT & 1S LiPo via CN1 & 3.0--4.2 V typical & nPM1300 VBAT. Battery pack must include PCM/protection.\\
VSYS & nPM1300 power path & USB or battery dependent & System input rail. Do not drive VSYS from an external source.\\
+1V8 & nPM1300 BUCK1 & 1.8 V & nRF52840 core/I/O domain and low-voltage logic as drawn.\\
+3.3V & nPM1300 BUCK2 & 3.3 V target & Application 3.3 V loads. VSET2 resistor must select 3.3 V at power-up.\\
+12V & MT3608 boost & 12 V & Solenoid only. Pulse load; verify diode/MOSFET stress and peak current.\\
+15V & MT3608 boost & 15 V & Piezo H-bridge only. Produces 30 Vpp differential drive.\\
VBUS & USB-C & 5 V nominal & nPM1300 VBUS input. Add local VBUS decoupling and USB ESD protection.\\
\bottomrule
\end{tabularx}
\end{table}

\subsection{nPM1300 Notes}

\begin{itemize}
  \item nPM1300 VBUS is the charger/system input. VBUSOUT is not required for charging.
  \item VBUSOUT is only a protected VBUS sense output for a host USB VBUS detect input. Leave NC if unused.
  \item Battery charging is firmware-enabled over TWI/I2C; configure charge current, termination voltage, NTC type, and VBUS current limit before production use.
  \item BUCK outputs must not be tied together.
  \item nPM1300 recommended BUCK external components from datasheet: 2.2 uH inductor, effective output capacitance at least 4 uF, low ESR.
  \item VSET pins are sampled at power-on only. Do not leave VSET1/VSET2 floating.
\end{itemize}

\subsection{BUCK Startup Configuration}

\begin{table}[h]
\small
\begin{tabularx}{\textwidth}{lllL}
\toprule
\textbf{Pin} & \textbf{Target} & \textbf{Required resistor to GND} & \textbf{Review note}\\
\midrule
VSET1 & 1.8 V & 47 k$\Omega$ & Verify against schematic/BOM.\\
VSET2 & 3.3 V & 250--500 k$\Omega$ & Use 330 k$\Omega$ nominal unless BOM already has another value in range.\\
\bottomrule
\end{tabularx}
\end{table}

\section{Battery and NTC}

\begin{table}[h]
\small
\begin{tabularx}{\textwidth}{lll}
\toprule
\textbf{CN1 pin} & \textbf{Net} & \textbf{Function}\\
\midrule
1 & VBAT & Battery positive\\
2 & NTC & 10 k$\Omega$ NTC to battery negative\\
3 & GND & Battery negative / system ground\\
\bottomrule
\end{tabularx}
\end{table}

CN1 is intended to be a 3-pin, 2.54 mm battery connector. Do not label this as JST-PH unless the selected part is actually PH-series; JST-PH is 2.0 mm pitch. If a 2.54 mm connector is used, identify the exact family/part number in the schematic and BOM. Silkscreen should identify the battery pinout as \texttt{+ NTC -}. Confirm footprint pin order against the mating cable before ordering.

\textbf{NTC requirements:}
\begin{itemize}
  \item Preferred pack: 1S LiPo, approximately 1000 mAh, protected, with 10 k$\Omega$ NTC.
  \item nPM1300 default temperature thresholds are intended for a 10 k$\Omega$ NTC profile.
  \item If no NTC is used, the nPM1300 NTC function must be disabled in firmware and the NTC pin tied as specified by the datasheet. For this product, using a real battery NTC is preferred.
\end{itemize}

\section{USB-C Interface}

\begin{table}[h]
\small
\begin{tabularx}{\textwidth}{lllL}
\toprule
\textbf{USB-C pins} & \textbf{Net} & \textbf{Destination} & \textbf{Notes}\\
\midrule
A4/A9/B4/B9 & VBUS & nPM1300 VBUS & Tie all VBUS contacts together. Local VBUS decoupling required.\\
A1/B12/B1/A12 & GND & Board GND & Tie all ground contacts together.\\
A5 & CC1 & nPM1300 CC1 & Connect directly; nPM1300 provides Rd/current detection.\\
B5 & CC2 & nPM1300 CC2 & Connect directly; nPM1300 provides Rd/current detection.\\
A6/B6 & USB\_D+ & nRF52840 USB D+ & Tie connector D+ contacts together. Use ESD protection.\\
A7/B7 & USB\_D- & nRF52840 USB D- & Tie connector D- contacts together. Use ESD protection.\\
A8/B8 & SBU1/SBU2 & NC & Unused for USB2 charging/data.\\
13/14 & Shield/EH & Chassis/GND strategy & Connector shell pads. Prefer direct GND or defined shield network near connector.\\
\bottomrule
\end{tabularx}
\end{table}

\textbf{Important:} Do not add external 5.1 k$\Omega$ CC pulldowns when CC1/CC2 are connected to the nPM1300; the nPM1300 implements the Type-C sink detection path internally.

\section{User Button, Ship Mode, and Reset}

\begin{table}[h]
\small
\begin{tabularx}{\textwidth}{llL}
\toprule
\textbf{Signal} & \textbf{Connection} & \textbf{Notes}\\
\midrule
SHPHLD & User button to GND & Wakes nPM1300 from Ship/Hibernate. Internal pull-up to VBAT/VBUS.\\
NPM\_INT & nPM1300 GPIO1 to nRF52840 & Used for PMIC events, including SHPHLD events if configured.\\
NPM\_RESET & nRF52840 to nPM1300 GPIO0 & MCU-controlled reset/control path. Keep separate from the user button.\\
\bottomrule
\end{tabularx}
\end{table}

Recommended user-button behavior:
\begin{itemize}
  \item Off/Ship: button pulls SHPHLD low and wakes the nPM1300.
  \item On: nPM1300 reports button event to nRF52840 through PMIC interrupt and TWI status registers.
  \item Firmware long-press action: nRF52840 commands nPM1300 to enter Ship mode.
  \item nPM1300 long SHPHLD press reset is enabled by default; decide whether firmware should disable or keep it.
\end{itemize}

\section{High-Current / High-Voltage Loads}

\begin{table}[h]
\small
\begin{tabularx}{\textwidth}{llllL}
\toprule
\textbf{Load} & \textbf{Rail} & \textbf{Driver} & \textbf{Control} & \textbf{Review notes}\\
\midrule
Solenoid lock & +12V & 2N7002 low-side & nRF52840 GPIO & Verify solenoid current, MOSFET pulsed SOA, flyback path, and pulse duration.\\
Piezo siren & +15V & MOSFET H-bridge & 2.9 kHz complementary PWM & Verify no shoot-through, gate resistor values, and piezo voltage limited to 30 Vpp.\\
Strobe LED & +3.3V & MOSFET low-side & PWM/strobe GPIO & Verify current limit topology, peak current, thermal rise, and BUCK2 current budget.\\
\bottomrule
\end{tabularx}
\end{table}

\section{Schematic Cross-Reference}

\begin{table}[h]
\small
\begin{tabularx}{\textwidth}{llL}
\toprule
\textbf{Item} & \textbf{Current schematic evidence} & \textbf{Review action}\\
\midrule
USB-C connector & USB1, \partno{C2765186}, 16-pin USB2 receptacle & Confirm shield pins 13/14 are passive and connected intentionally.\\
USB CC routing & CC1/CC2 labels present at USB1 and nPM1300 & Confirm no external CC Rd resistors are also fitted.\\
USB data & USB\_D+ / USB\_D- labels present at USB1 and nRF52840 & Confirm D+/D- polarity and ESD part orientation.\\
Battery connector & CN1 intended as 2.54 mm 3-pin & Confirm exact connector family/part number and pin order 1=VBAT, 2=NTC, 3=GND against footprint/cable.\\
Battery section text & Schematic text says ``Battery JST 2.54mm 1x3P'' & Keep 2.54 mm if that is the selected pitch, but remove ``JST-PH'' wording unless using a real PH-series 2.0 mm connector.\\
VSET2 & VSET2 net present on nPM1300 & Confirm resistor value is 250--500 k$\Omega$ for 3.3 V startup.\\
VBUSOUT & nPM1300 VBUSOUT pin present & If NC, no capacitor required; if used for host VBUS sense, add required decoupling.\\
SHPHLD & SW1 near nPM1300; SHPHLD pin present & Confirm SW1 pulls SHPHLD to GND only, not directly to nRF GPIO at VBAT voltage.\\
PMIC interrupt & NPM\_INT label between nPM1300 and nRF52840 & Confirm firmware will enable SHPHLD/button events if button has no direct MCU GPIO.\\
\bottomrule
\end{tabularx}
\end{table}

\section{Open Review Items}

\begin{enumerate}
  \item Confirm nPM1300 package choice and assembly capability. QFN is preferred for prototype inspectability; WLCSP is acceptable only with proven assembler capability.
  \item Confirm all nPM1300 required decoupling capacitors and BUCK passives match the datasheet/reference design.
  \item Resolve battery connector naming: 2.54 mm is acceptable if it matches the cable, but it should not be called JST-PH unless it is a real 2.0 mm PH-series part.
  \item Confirm USB-C shield grounding strategy and ESD placement before layout.
  \item Confirm charge profile for the selected cell: chemistry, VTERM, ICHG, JEITA thresholds, and NTC beta selection.
  \item Confirm +3.3V BUCK2 load budget. nPM1300 BUCK outputs are 200 mA class rails; strobe LEDs may need a separate driver/rail if peak current exceeds budget.
  \item Confirm +12V solenoid and +15V siren boost converters are not enabled during low battery, brownout, or USB current-limit states without firmware gating.
  \item Confirm all nets driven from VBAT/VBUS into nRF52840 pins are level-safe. SHPHLD is not safe to tie directly to a 3.3 V-only GPIO.
  \item Run KiCad ERC after connector pin types are cleaned: connector pins passive, USB shell passive, explicit PWR\_FLAGs only where needed.
\end{enumerate}

\end{document}
